\documentclass[11pt,letterpaper]{article}
\usepackage[margin=1in]{geometry}
\usepackage{amsmath,amssymb,mathtools}
\usepackage{microtype}
\usepackage{enumitem}
\usepackage{xcolor}
\usepackage{hyperref}
\hypersetup{colorlinks=true,linkcolor=blue,urlcolor=blue}
\setlength{\parskip}{0.6em}
\setlength{\parindent}{0pt}
\newcommand{\vect}[1]{\boldsymbol{#1}}
\newcommand{\Smat}{\mathbf{S}}
\newcommand{\fvec}{\mathbf{f}}
\newcommand{\yvec}{\mathbf{y}}
\newcommand{\evec}{\mathbf{e}}
\newcommand{\Rmat}{\mathbf{R}}
\newcommand{\pvec}{\mathbf{p}}
\begin{document}
\begin{center}
{\Large\bfseries Least-Squares Filter Design}\\[4pt]
{\large ERTH 467 --- Geoscientific Signal Processing}\\[3pt]
\emph{From prediction error to the normal equations}
\end{center}

\section{What problem are we trying to solve?}
Suppose we know an input signal $s$ and a desired output (or \emph{target}) signal $y$.
We seek the coefficients of a linear filter $f$ such that convolving $s$ with $f$
produces an output as close as possible to $y$:
\begin{equation}
    \widehat y_n = \sum_{k=0}^{M-1}f_k\,s_{n-k}.
    \label{eq:convolution}
\end{equation}
Here $M$ is the number of unknown filter coefficients. We collect samples into vectors,
and write the convolution as a matrix multiplication:
\begin{equation}
    \boxed{\widehat{\yvec}=\Smat\fvec.}
\end{equation}
For $N$ output samples, $\Smat$ is an $N\times M$ matrix, $\fvec$ is an $M\times 1$
vector, and $\yvec$ and $\widehat{\yvec}$ are $N\times 1$ vectors.
Each column of $\Smat$ is a shifted copy of the input signal $s$.

\textbf{Matrix convention.} For finite, zero-padded, causal convolution, $\Smat$ has
Toeplitz structure, for example
\begin{equation}
\Smat=\begin{bmatrix}
s_0&0&0\\
s_1&s_0&0\\
s_2&s_1&s_0\\
s_3&s_2&s_1\\
\vdots&\vdots&\vdots
\end{bmatrix}.
\end{equation}
This is the type of matrix displayed in the lecture slide. If we instead assume
\emph{periodic} indexing, $s_{n-k}$ is interpreted modulo $N$ and $\Smat$
can be a \emph{circulant} matrix. For instance, with $N=M=3$,
\begin{equation}
\Smat_{\mathrm{circ}}=\begin{bmatrix}
s_0&s_2&s_1\\
s_1&s_0&s_2\\
s_2&s_1&s_0
\end{bmatrix}.
\end{equation}
\textbf{The least-squares derivation below is identical in either case.}

\section{Define the prediction error}
For each sample, define the prediction error as the difference between the
observed target and the filtered signal:
\begin{equation}
    e_n=y_n-\widehat y_n=y_n-\sum_{k=0}^{M-1}S_{nk}f_k.
\end{equation}
In vector form,
\begin{equation}
    \boxed{\evec=\yvec-\Smat\fvec.}
\end{equation}
There are generally more data samples than unknown filter coefficients ($N\ge M$),
so it may not be possible to make every $e_n$ exactly zero.
Instead, choose $\fvec$ to minimize the \emph{sum of squared prediction errors}:
\begin{equation}
    \boxed{J(\fvec)=\sum_{n=0}^{N-1}e_n^2
    =\evec^T\evec
    =(\yvec-\Smat\fvec)^T(\yvec-\Smat\fvec).}
    \label{eq:cost}
\end{equation}
This is the least-squares objective. It is a \emph{scalar} function of the
$M$ unknown coefficients $f_0,\ldots,f_{M-1}$.

\section{Find the minimum by differentiating}
The key idea is the same as in elementary calculus: at a differentiable minimum,
the derivative with respect to each adjustable coefficient is zero.
With several unknowns, we collect all partial derivatives in the \emph{gradient}:
\begin{equation}
    \nabla_{\fvec}J=
    \begin{bmatrix}
    \partial J/\partial f_0\\
    \partial J/\partial f_1\\
    \vdots\\
    \partial J/\partial f_{M-1}
    \end{bmatrix},\qquad
    \text{and require}\qquad \nabla_{\fvec}J=\mathbf{0}.
\end{equation}

\subsection*{Step 1: Expand the squared error}
Using $(\mathbf{a}-\mathbf{b})^T=\mathbf{a}^T-\mathbf{b}^T$,
\begin{align}
J(\fvec)
&=(\yvec^T-\fvec^T\Smat^T)(\yvec-\Smat\fvec)\\
&=\yvec^T\yvec-\yvec^T\Smat\fvec
-\fvec^T\Smat^T\yvec+\fvec^T\Smat^T\Smat\fvec.
\end{align}
The two middle terms are both \emph{scalars} and equal because one is the
transpose of the other. Therefore,
\begin{equation}
    \boxed{J(\fvec)=\yvec^T\yvec
    -2\fvec^T\Smat^T\yvec
    +\fvec^T\Smat^T\Smat\fvec.}
    \label{eq:expanded}
\end{equation}

\subsection*{Step 2: Differentiate with respect to the filter}
Treat $\Smat$ and $\yvec$ as known; only $\fvec$ is variable.
The three terms of Eq.~\eqref{eq:expanded} have gradients
\begin{align}
\nabla_{\fvec}(\yvec^T\yvec)&=\mathbf{0},
&&\text{(independent of $\fvec$)},\\
\nabla_{\fvec}(-2\fvec^T\Smat^T\yvec)&=-2\Smat^T\yvec,
&&\text{(linear in $\fvec$)},\\
\nabla_{\fvec}(\fvec^T\Smat^T\Smat\fvec)&=2\Smat^T\Smat\fvec,
&&\text{(quadratic in $\fvec$)}.
\end{align}
The last identity holds because $\Smat^T\Smat$ is symmetric.
Combining the three derivatives gives
\begin{equation}
    \boxed{\nabla_{\fvec}J
    =-2\Smat^T\yvec+2\Smat^T\Smat\fvec
    =-2\Smat^T(\yvec-\Smat\fvec).}
    \label{eq:gradient}
\end{equation}
\textbf{Interpretation:} the derivative tells us how the total squared error
changes when we slightly modify each coefficient of the filter.

\subsection*{Step 3: Set the gradient to zero}
At a minimum,
\begin{align}
\nabla_{\fvec}J&=\mathbf{0},\\
-2\Smat^T\yvec+2\Smat^T\Smat\fvec&=\mathbf{0},\\
\boxed{\Smat^T\Smat\fvec=\Smat^T\yvec.}
\label{eq:normal}
\end{align}
These are the \textbf{normal equations}, a system of $M$ linear equations
for the $M$ unknown filter coefficients.

If $\Smat^T\Smat$ is invertible, the formal solution is
\begin{equation}
    \boxed{\fvec_{\mathrm{LS}}=(\Smat^T\Smat)^{-1}\Smat^T\yvec.}
    \label{eq:solution}
\end{equation}
This is the expression on the lecture slide. In an actual numerical program,
we generally \emph{solve} the linear system rather than explicitly constructing
an inverse. If $\Smat$ has linearly dependent columns, the inverse does not exist;
a pseudoinverse or regularization can be used instead.

\section{Why does this stationary point minimize the error?}
The curvature of the objective is described by its Hessian:
\begin{equation}
    \nabla_{\fvec}^2J=2\Smat^T\Smat.
\end{equation}
For any vector $\mathbf{v}$,
\begin{equation}
    \mathbf{v}^T\Smat^T\Smat\mathbf{v}
    =(\Smat\mathbf{v})^T(\Smat\mathbf{v})
    =\|\Smat\mathbf{v}\|_2^2\ge 0.
\end{equation}
Thus $J$ is a convex quadratic function, and a solution of the normal equations
is a \emph{global minimum}, not a maximum. The minimum is unique when the columns
of $\Smat$ are linearly independent.

There is also a useful geometric interpretation. From Eq.~\eqref{eq:gradient},
at the optimum
\begin{equation}
    \boxed{\Smat^T\evec_{\mathrm{LS}}=\mathbf{0}.}
\end{equation}
The remaining prediction error is orthogonal to every column of $\Smat$.
In other words, no additional linear combination of the available shifted
signal samples can reduce the squared error.

\section{Connection with correlation}
It is common to introduce
\begin{equation}
    \Rmat=\Smat^T\Smat,
    \qquad \pvec=\Smat^T\yvec.
\end{equation}
Then the normal equations become
\begin{equation}
    \boxed{\Rmat\fvec=\pvec.}
\end{equation}
For a convolution matrix, entries of $\Rmat$ contain inner products of shifted
copies of the input signal (related to \emph{autocorrelation}), while entries of
$\pvec$ contain inner products between shifted input samples and the target
(related to \emph{cross-correlation}). Exact edge behavior depends on whether
we use finite zero-padded or periodic convolution.

\section{Optional: check the derivative one coefficient at a time}
For students who prefer ordinary scalar differentiation, start directly from
\begin{equation}
J=\sum_{n=0}^{N-1}\left(y_n-\sum_{k=0}^{M-1}S_{nk}f_k\right)^2.
\end{equation}
Differentiate with respect to one coefficient $f_j$ using the chain rule:
\begin{align}
\frac{\partial J}{\partial f_j}
&=2\sum_{n=0}^{N-1}
\left(y_n-\sum_{k=0}^{M-1}S_{nk}f_k\right)(-S_{nj})\\
&=-2\sum_{n=0}^{N-1}S_{nj}e_n.
\end{align}
Setting this equal to zero for every $j=0,\ldots,M-1$ gives
\begin{equation}
\sum_{n=0}^{N-1}S_{nj}\sum_{k=0}^{M-1}S_{nk}f_k
=\sum_{n=0}^{N-1}S_{nj}y_n,
\end{equation}
which is precisely the $j$th row of $\Smat^T\Smat\fvec=\Smat^T\yvec$.
The scalar derivation and the matrix derivation are therefore the same calculation.

\section*{Main result to remember}
A least-squares filter chooses the coefficients $\fvec$ that make the
predicted signal $\Smat\fvec$ as close as possible to the target $\yvec$
in the \emph{sum-of-squared-errors} sense:
\begin{equation*}
\underbrace{J(\fvec)=(\yvec-\Smat\fvec)^T(\yvec-\Smat\fvec)}_{\text{define error}}
\quad\xrightarrow{\;\nabla_{\fvec}J=0\;}\quad
\underbrace{\Smat^T\Smat\fvec=\Smat^T\yvec}_{\text{solve for filter}}.
\end{equation*}
\end{document}
